\section{The direct algebra backend and proof exchange}

\subsection{What the checked library now provides}

The new \path{Logic/PrimitiveRecursiveArithmetic/Derivations} library changes
the practical starting point for algebraic formalization.  It is not another
research packet: its README states that every listed theorem has a concrete
closed type \([]\vdash_{\PRAAx}\varphi\), is checked by Lean's kernel, and has no
\path{sorry}.  Its reusable layers include:

\begin{itemize}
\item PR symbols and defining equations for \(+\), \(\cdot\), predecessor, and
truncated subtraction;
\item arbitrary-term congruence;
\item open quantifier-free induction \path{ind};
\item instantiation of closed derivations \path{inst};
\item commutative-semiring laws and additive cancellation;
\item a quotient \(Q\Gamma\) of terms by derivable equality;
\item integer terms as differences and the ring \(\mathsf{Groth}(Q\Gamma)\);
\item \path{ipra_ring}, which turns a ring-normalization proof in the
metatheory into an object-level derivation.
\end{itemize}

The first checked applications are polynomial identities behind Heron's formula
and the law of cosines.  The method is more general than those examples.

\subsection{Why metatheoretic ring normalization is legitimate}

The tactic does not add Mathlib's ring theorem as an axiom of iPRA.  The object
level already proves the required semiring and cancellation laws.  Terms modulo
provable equality form a cancellative commutative semiring.  Integer pairs map
injectively into its Grothendieck ring.  Mathlib's \path{ring} tactic proves an
equality in that metatheoretic ring; the library then uses the proved
injectivity and quotient soundness to construct the corresponding
\(\Prf\;\PRAAx\) derivation.

Thus the tactic is a proof-term generator.  Lean checks the generated object
proof.  This is exactly the right use of a powerful host tactic.

\subsection{Polynomial identity node}

A polynomial identity node has:

\begin{enumerate}
\item a finite list of natural or integer variables;
\item two expression trees using addition, multiplication, numerals, and
integer pair operations;
\item the equation formula between their compiled \(\PTm\) terms;
\item a direct \(\Prf\;\PRAAx\) derivation produced by \path{ipra_ring}.
\end{enumerate}

Theorem~\ref{thm:import} imports the result into the native graph.  Side
conditions such as denominator positivity, divisibility, nonzeroness, or order
are not ring identities.  They remain separate clauses and are combined by the
graph compiler.

\subsection{A proof-exchange normal form}

To make direct and native developments interoperable, use the following normal
form for a local lemma.

\begin{definition}[Exchange package]
An exchange package consists of an ordinary formula \(\varphi\), its native
translation \(A=\NFm.\mathsf{unQ}(\NFm.\mathsf{ofPRA}(\varphi))\), and one of:
\begin{enumerate}
\item a direct derivation \(h:\Gamma\vdash_{\PRAAx}\varphi\), to be imported by
Theorem~\ref{thm:import}; or
\item a native derivation \(k:\KPrf(\Delta,A)\), to be exported by
Theorem~\ref{thm:export}.
\end{enumerate}
The package records the round-trip equivalence supplied by
\path{pra_native_roundTrip} or the native round-trip theorem.
\end{definition}

Application libraries may expose both forms, but only one proof need be
maintained.  The compiler can choose whichever form is more convenient.

\subsection{Inequalities and order}

Polynomial equality is only one part of quantitative mathematics.  Natural
order may be encoded by truncated subtraction or by an existential difference:
\[
   a\le b\quad\Longleftrightarrow\quad a\mathbin{\dot-}b=0,
   \qquad\text{or}\qquad \exists c\;a+c=b.
\]
For bounded calculations, the first is convenient and decidable.  For algebraic
reasoning, the second may compose better with addition.

A practical next layer for \path{IPRADerivations} should include:

\begin{itemize}
\item reflexivity, transitivity, antisymmetry, and totality of natural order;
\item monotonicity of addition and multiplication by nonnegative terms;
\item strict order and positive-product lemmas;
\item clearing denominators under positivity guards;
\item square nonnegativity and sum-of-squares arguments;
\item bounded linear arithmetic proof generation.
\end{itemize}

Until that layer exists, the present framework treats these as explicit local
node obligations.  Many can be proved by quantifier-free induction using
\path{Core/Basic} and \path{Core/Semiring}.

\subsection{Finite sums and products}

A large class of identities uses bounded sums and products.  Define PR folds
\[
 \Sigma_f(0)=0,\quad \Sigma_f(n+1)=\Sigma_f(n)+f(n),
\]
\[
 \Pi_f(0)=1,\quad \Pi_f(n+1)=\Pi_f(n)f(n).
\]
The needed library lemmas include splitting an interval, reindexing by a bounded
bijection, distributivity over sums, telescoping, and finite double-sum swaps.
Each is proved by quantifier-free induction.

Once these lemmas are available, the ring backend can discharge the local
polynomial step while induction handles the bound.  This division mirrors
ordinary mathematical practice: induction establishes the finite combinatorial
shape, and ring normalization establishes each algebraic identity.

\subsection{Fixed-dimensional determinant example}

For a fixed dimension \(d\), define the determinant by the Leibniz formula
\[
  \det A=\sum_{\sigma\in S_d}\operatorname{sgn}(\sigma)
           \prod_{i<d}a_{i,\sigma(i)}.
\]
The permutation list and all products have constant finite size.  Determinant
multilinearity, alternation, and Laplace expansion become finite polynomial
identities.  They can be proved directly by \path{ipra_ring} after the finite
sum has been expanded, or modularly using a bounded-sum library.

For a concrete \(d\), code generation can expand the entire identity and let
the ring tactic produce the proof.  For variable \(d\), a generic permutation
and sign library is preferable.  The native graph compiler handles the finite
loop and table dependencies in either approach.

\subsection{Certificate checking versus symbolic derivation}

Two complementary strategies are useful.

\paragraph{Symbolic derivation.}
Prove the identity for variables.  This gives a small reusable theorem and is
ideal for polynomial algebra.

\paragraph{Certificate checking.}
For a concrete input, provide a witness such as an inverse matrix, factorization,
Smith form, root-count chain, or Groebner reduction and verify it by finite
arithmetic.  This often avoids formalizing the algorithm that found the
certificate.

A general framework should support both.  The producer may be sophisticated or
external; the verifier and finalizer remain object-level.  A symbolic theorem
can also prove the verifier sound once for all certificates.

\subsection{Recommended extension modules}

The direct derivation library would gain the most leverage from the following
next modules, in dependency order.

\begin{enumerate}
\item \path{Core/Order}: natural order, positivity, monotonicity, and denominator
clearing.
\item \path{Core/Pairing}: a fixed pair/list code with projection equations and
bounded table operations.
\item \path{Core/FiniteFold}: sums, products, bounded Boolean folds, and
reindexing.
\item \path{Core/Rationals}: extensional rational arithmetic and order.
\item \path{Core/Matrix}: fixed and variable finite matrices, dot products, and
matrix identities.
\item \path{Core/NativeExchange}: the import/export lemmas of this report.
\item \path{Core/Trace}: a reusable object-level clocked-trace theorem.
\end{enumerate}

These modules would turn many proof packets from isolated blueprints into
applications of a shared checked stack.
